\documentclass[12pt,a4paper]{ctexart}
\usepackage{amsmath,amssymb}
\usepackage{booktabs}
\usepackage{array}
\usepackage{geometry}
\usepackage{enumitem}
\usepackage[hidelinks]{hyperref}
\geometry{margin=2.5cm}
\setlength{\parskip}{0.5em}
\setlength{\parindent}{2em}
\linespread{1.25}
\renewcommand{\arraystretch}{1.15}

\ctexset{
  section = {format = {\heiti \zihao{3}}},
  subsection = {format = {\heiti \zihao{4}}},
  subsubsection = {format = {\heiti \zihao{-4}}},
}

\newcommand{\ok}{$\checkmark$}
\newcommand{\warn}{$\triangle$}
\newcommand{\fix}{$\times$}
\newcommand{\dg}{\ensuremath{^\circ}}

\begin{document}

\begin{center}
{\heiti\zihao{3} A 题论文验证报告\par}
\vspace{0.4em}
{\zihao{5} 数据核对 \,\textbullet\, 数值复现 \,\textbullet\, 独立交叉验证 \,\textbullet\, 问题清单（2026-09-12）\par}
\end{center}

\section{一、验证范围与方法}

本报告对《A题论文》(A题论文.tex, 2026-09-12 版) 的\textbf{数据引用、数值结果与物理量级论断}逐项复核,采用四条独立证据链:

\begin{enumerate}[label=(\arabic*),leftmargin=2.5em]
\item \textbf{附件数据核对}:附件 1、附件 2 原始 xlsx 直接读取,与论文引用的每个数值对照;
\item \textbf{求解器复现}:q1--q4 四个求解器全部重新运行,输出与原日志逐行比对,再与 result1--4.xlsx 及论文表 1--6 三方互查;
\item \textbf{独立交叉验证}:Bessel 级数解析解用 scipy 独立复算(与论文 7.1 节使用不同代码);离散守恒、收敛比、渐近行为按论文 7.2--7.4 节的方法复核;
\item \textbf{物理量级复核}:7.6 节汽化潜热分析的每个数字独立重算(蒸汽压、湿球温度、Arrhenius 因子、热量比)。
\end{enumerate}

\section{二、附件数据核对（全部通过）}

\begin{center}
{\zihao{5}\begin{tabular}{lll}
\toprule
论文引用 & 附件实际数据 & 结论 \\
\midrule
附件 1 共 241 点、0--14400 s & 242 行(含表头)、每 60 s 一点 \ok & 一致 \\
$T_a$(1800 s)$=41.513\dg$C & $41.513\dg$C \ok & 一致 \\
$T_a$(14400 s)$=50.165\dg$C, $C_a=0.04986$ & 50.165 / 0.04986 \ok & 一致 \\
前 1800 s 烘房 28$\to$41.513\dg C & 28$\to$41.513\dg C \ok & 一致 \\
附件 2 共 145 点、每 1800 s & 146 行(含表头)、0--259200 s \ok & 一致 \\
30 min 收缩约 6.4\% & $R$(1800 s)$=1.873$ cm, 收缩 6.35\% \ok & 一致 \\
全程收缩约 40.1\% & 末值 1.198 cm, 收缩 40.10\% \ok & 一致 \\
2 cm$\to$约 1.5 cm 在前 3.5 h & $R$(12600 s)$=1.51$ cm \ok & 一致 \\
表 6 中 6 h 起 1.5/2 cm 列空缺 & $R$(21600 s)$=1.374$ cm \ok & 正确 \\
``72 h 后稳定于 1.198 cm'' & \textbf{67 h}(241200 s)起即恒定 \warn & 措辞宜改 \\
\bottomrule
\end{tabular}}
\end{center}

另核对 7.5 节文献互证数据:丹参 $D_{\mathrm{eff}}=1.7\times10^{-8}\sim9.4\times10^{-8}$ m$^2$/s、$E_a=40.0$ kJ/mol;鲜地黄 $E_a=34.94$ kJ/mol;白术 $E_a=37.47$ kJ/mol,与文献清单记载完全一致 \ok。

\textbf{重要更正}:仓库旧版学习材料中有两处错误数据,本次已修正:①《A题物理化学补习指南》中``前 1800 s 烘房仅升至约 29.7\dg C''是错的(实为 $41.513\dg$C;$29.7\dg$C 只是约 $t\approx190$ s 时的值);②《A题方法详解》中``$t=72000$ s 时 $R=1.198$ cm''是错的(实为 1.210 cm,1.198 cm 要到 67 h 才达到)。论文本身未采用这两处错误数据。

\section{三、数值结果验证（全部通过）}

\subsection{3.1 求解器复现}

四问求解器重跑结果与原日志\textbf{逐行一致}:
Q1 仅附件 2 打印行因本次修复标签而不同(数值不变);
Q2、Q3 完全一致;
Q4 仅末行标签由 ``54.0h'' 改为 ``50.7h''(实际时间本来就是 $t^*$)。
四个 result xlsx 由重跑重新生成,数值不变。

\subsection{3.2 论文表格与 result 文件三方一致}

\begin{center}
{\zihao{5}\begin{tabular}{llll}
\toprule
表格 & result 文件 & 抽查点 & 结论 \\
\midrule
表 1/表 2(问题 1) & result1.xlsx & 100、300、1800 s 全部 5 列 + 日志 7 行全比对 & 一致 \ok \\
表 3/表 4(问题 2) & result2.xlsx & 0.5、1、3 h 全部 5 列 + 日志 6 行全比对 & 一致 \ok \\
表 5(问题 3) & result3.xlsx & 6、12、54 h + 烘干结束行 & 一致 \ok \\
表 6(问题 4) & result4.xlsx & 6、48 h + 烘干结束行 + 表面列 & 一致 \ok \\
\bottomrule
\end{tabular}}
\end{center}

result 文件与附件 3 模板格式一致(result1/2 含``温度''``水分浓度''两表,时间从 1 s 起;result3/4 单表,时间从 60 s 起——模板本身即从这些时间开始,无缺失)。

\subsection{3.3 独立交叉验证}

\textbf{(1) Bessel 解析解独立复算。}对问题 1 热方程,本报告用 scipy 独立求解特征方程 $\beta J_1(\beta)=\mathrm{Bi}\,J_0(\beta)$ 的前 40 个正根并直接求和级数(与论文 7.1 节完全独立的代码),得 $t=1800$ s 时中心 $T=37.8050\dg$C、表面 $T=39.4483\dg$C,与求解器数值解\textbf{四位小数完全一致}。这独立确认了离散格式、圆心处理与 Robin 表面行的正确性。

\textbf{(2) 离散质量守恒。}重跑复现:问题 1 闭合差 $1.5\times10^{-18}$、问题 3 为 $1.7\times10^{-17}$(机器精度);问题 4 为 $4.2\times10^{-7}$(相对总失水量 $8\times10^{-4}$,源于动边界半步中点的 $O(\Delta t^2)$ 配对误差)。与论文 7.3 节数字一致 \ok。

\textbf{(3) 收敛性。}重跑复现:问题 1 表 2 收敛比 4.3(空间二阶);问题 2 表 4 收敛比 4.0;问题 4 $\zeta$ 网格加密最大差 $5.33\times10^{-4}$;各时间收敛差 $10^{-5}\sim10^{-6}$ 量级。与论文 7.2 节数字一致 \ok。

\textbf{(4) 渐近行为。}常数边界 $(50.165\dg\mathrm{C},0.04986)$ 下长时程:中心温度四位小数收敛到 50.1650;表面水分浓度由 $10^5$ s 的 0.0616 缓慢趋近 0.04986($3\times10^5$ s 为 0.0512),与 $D\to0$ 时呈指数拖尾的预期一致 \ok。

\textbf{(5) 烘干判据实现。}问题 3 判据 ``各处 $C\le0.15$'' 用\textbf{未四舍五入的原始值}判断(代码 `if C.max() <= CTARGET`),未受四舍五入污染 \ok。注意:result3.xlsx 中按四舍五入数据看,204960 s 时中心已显示 0.1500,但精确穿越发生在 205082.5 s——判据实现正确,四舍五入只是显示假象。

\section{四、物理量级论断复核（7.6 节等）}

\begin{center}
{\zihao{5}\begin{tabular}{p{5.2cm}p{2.6cm}p{2.6cm}p{2.4cm}}
\toprule
论断(论文) & 论文数值 & 独立重算 & 结论 \\
\midrule
活化能 $E_a=3850R_g$ & 32.0 kJ/mol & 32.0 kJ/mol & 一致 \ok \\
50\dg C 饱和含湿量 & 0.0868 & 0.0863 & 差 0.6\%,可忽略 \ok \\
相对湿度 RH & 约 57\% & 57.8\% & 一致 \ok \\
湿球温度 $T_{wb}$ & 约 $40\dg$C & \textbf{41.55}\dg C & 量级可接受 \warn \\
蒸发吸热/对流显热 & \textbf{约 6} & \textbf{4.95}(论文自己的数);中期干基密度 $\rho_d=\rho/(1+C)\approx390\sim530$ 时为 3--4 & 偏高,建议改``约 3--5'' \warn \\
Arrhenius 因子 $T$:50$\to$40\dg C & 0.68 & 0.6794 & 一致 \ok \\
时长敏感性 & \textbf{约 +30\%}(74 h) & 按 $t^*\propto1/D$ 标度为 \textbf{+47\%}($\approx$84 h$\approx$3.5 天) & 偏乐观,见问题 2 \fix \\
附录 3 的 $D$ 比附录 2 ``大约一个量级'' & 约 $10^{-8}$ vs $4\times10^{-9}$ & 比值 2.7--3.6 倍($C$=1--2.55) & 应改``约 2--3 倍'' \warn \\
附录 4 的 $D$ 比附录 3 小 3--5 倍 & 3--5 倍 & 2.1--5.4 倍(随 $C$) & 基本准确 \ok \\
收缩使 $D/R^2$ 增大 & 2.8 倍 & $(2/1.198)^2=2.787$ & 一致 \ok \\
$\alpha=k/(\rho c_p)$、$R^2/\alpha$ & $\sim2\times10^{-7}$、$\sim2000$ s & $1.45\times10^{-7}$、2762 s & 量级一致,略偏 \warn \\
$\sqrt{Dt}\approx2.6$ mm & 2.6 mm & 2.6 mm($D=3.9\times10^{-9}$) & 一致 \ok \\
边界层 $\delta\sim D/h_m\approx5$ mm & 5 mm & 4.9--6.1 mm & 一致 \ok \\
$t^*$ 缩短 6.3 h(11\%) & 6.3 h/11\% & 6.27 h/11.0\% & 一致 \ok \\
\bottomrule
\end{tabular}}
\end{center}

\section{五、发现的问题清单}

\subsection{5.1 需修正（影响结论表述）}

\textbf{问题 1(高):$t^*=205082$ s 应为 205082.5 s。}q3 求解器时间步 $\Delta t=2.5$ s,判据首次满足于第 82033 步,即 $t^*=82033\times2.5=205082.5$ s;求解器日志用 `\%d` 打印截断了小数位,论文沿用了截断值。而论文同时给出的 56.9674 h 恰好对应 205082.5 s($205082.5/3600=56.96736$)——论文内部不一致。建议将摘要与 6.3 节的 ``205082 s'' 改为 ``205082.5 s''(或若取整数秒则写 205083 s 并保留 56.9674 h)。Q4 的 $t^*=182495$ s 是 $\Delta t=5$ s 的整倍数,无此问题 \ok。

\textbf{问题 2(高):7.6(3) 的时长敏感性偏乐观,``不改变 2--3 天结论''表述过强。}若表面温度因汽化吸热由 50\dg C 降至 40\dg C,Arrhenius 因子减小 0.68(已确认),扩散系数下降 32\%;降速干燥段 $t^*\propto R^2/D$,时长应增加约 $1/0.68-1\approx47\%$,即 $56.97\to84$ h(约 3.5 天),而非论文所称 ``约 30\%(74 h)''。即便取论文自己的 74 h,也已达 3.08 天,略超题述``2--3 天''上限。建议改写为:敏感性约 30\%--50\%,显式计入潜热可能使时长增至 3 天以上;这恰恰强化了``题给经验式必须以有效参数吸收潜热效应''的论证,同时应坦承这一不确定度,而非断言``不改变结论次序''。

\subsection{5.2 建议修正（提升严谨性）}

\textbf{问题 3:}7.6 节 $q_{\mathrm{evap}}/q_{\mathrm{sens}}$ ``约 6'' 按论文自列数值代入实为 4.95;若按干燥中期的干基密度 $\rho_d=\rho(C)/(1+C)\approx390\sim530$ kg/m$^3$ 计算则为 3--4。建议改为``约 3--5(与对流显热同量级)''并注明密度取法。定性结论(潜热不可忽略、被有效参数吸收)不受影响。

\textbf{问题 4:}6.2 节``附录 3 的 $D$ 比附录 2 大约一个量级''——实际比值为 2.7--3.6 倍($C$=1--2.55 kg/kg 时),建议改为``大约 2--3 倍''(在低含水率端比值更大)。

\textbf{问题 5:}7.2 节未披露 $t^*$ 的网格敏感性:问题 3 在 0.5 mm 与 0.25 mm 网格上 $t^*$ 分别为 56.7069 h 与 56.9674 h,相差 0.26 h;按二阶收敛外推 $t^*\approx57.05$ h。报告 $t^*$ 到四位小时(56.9674 h,即 $\pm0.36$ s 精度)超出了网格不确定度。建议:7.2 节补充 $t^*$ 的网格收敛数据,并在摘要/6.3 将 $t^*$ 表述为``约 57.0 h''(或保留四位小数但注明网格不确定度约 $\pm0.1$ h)。Q4 同理。

\textbf{问题 6:}表 6 说明``$d>R(t)$ 处记为 ---''与 result4.xlsx 实际\textbf{留空}不一致;建议在 xlsx 中写入 ``---'' 或把论文措辞改为``留空''。另 result4.xlsx 末尾有一列多余空列(第 24 列),可清理。

\textbf{问题 7(措辞):}5.3 节与附录``72 h 后稳定于 1.198 cm''——数据实际自 67 h(241200 s)起恒定;建议改``67 h 起恒定于 1.198 cm(数据至 72 h)''。湿球温度``约 40\dg C''宜改``约 $41.5\dg$C'';$\alpha\sim2\times10^{-7}$ 宜改 $\sim1.5\times10^{-7}$(或保留,量级一致)。

\subsection{5.3 本次已修复}

\begin{enumerate}[label=(\arabic*),leftmargin=2.5em]
\item q1\_solver.py 附件 2 核对打印行误把末点值标为``t=72000s''(代码取 $d_2[-1,1]$,即 259200 s 的 1.198,标签却写 72000 s)——已修正标签,数值无影响;
\item 《A题方法详解》``$t=72000$ s 时 $R=1.198$ cm,此后恒定''——已改为``$t=72000$ s 时 $R=1.210$ cm,$t=241200$ s(67 h)起恒定于 1.198 cm'';
\item 《A题物理化学补习指南》``前 1800 s 烘房仅升至约 29.7\dg C''——已改为正确数据($41.513\dg$C),并注明 $29.7\dg$C 只是约 $t\approx190$ s 时的值;
\item 仓库学习类文件归档:学习路线图、方法详解、评估报告的 .tex 源文件已从 ProblemFolder 根目录移入 LearningMaterial,重复 PDF 与编译垃圾(aux/log/out)已删除。
\end{enumerate}

\section{六、总体结论}

\begin{enumerate}[label=(\arabic*),leftmargin=2.5em]
\item \textbf{模型与题设严格对应}:四问分别采用附录 2/3/4 参数与附件 1/2 数据,无混用;两阶段烘房条件拼接、判据实现、动边界随体坐标处理均正确。
\item \textbf{数值结果可信}:表 1--6 与 result 文件完全一致;四问求解器可复现;Bessel 解析解独立复算四位小数一致;离散守恒达机器精度;空间二阶收敛。问题 1、2 的全部输出值可以直接定稿。
\item \textbf{核心结论成立}:$t^*=56.97$ h(问题 3)与 50.69 h(问题 4)落在题述``2--3 天''范围内,收缩缩短烘干时长约 11\% 的结论由数据支撑。
\item \textbf{待修正}:(1) $t^*=205082\to205082.5$ s;(2) 7.6 节敏感性改 30\%--50\% 并软化``不改变 2--3 天''表述。其余问题 3--7 为严谨性改进,不影响结论。
\end{enumerate}

\end{document}
